\subsection{1. Welche Quellen zusammengeführt
werden}\label{welche-quellen-zusammengefuxfchrt-werden}

\begin{enumerate}
\def\labelenumi{\arabic{enumi}.}
\tightlist
\item
  Die Benutzeranlage
  \codeword{addbf22d-4764-4685-bdfe-f422bfc563d0/pasted-text.txt},
  eingefroren unter \codeword{sources/user\_attachment.txt}.
\item
  Die externe Untersuchung
  \codeword{Universalraum\_Beweisversuch\_2026-09-15.md}:
  Spektralzählung, Primzahl-Dynamik, Phasencocycle und kontrollierte
  Polnäherung.
\item
  Die externe Untersuchung
  \codeword{Universalraum\_Urspruenglicher\_Austausch\_2026-09-15.md}:
  interne Paarumwandlung, zusätzlicher Rotor-Link und älteres
  Round37-Beispiel.
\item
  Der abgesicherte native Stand v1.6.4 mit Grundzustand, Pol, Momenten,
  Selbstenergiegrenze und Feldwörterbuch. Das frühere vollständige
  Prüfpaket liegt unverändert bei. Seine komplette
  Konfigurationsenumeration wurde \textbf{nicht erneut} in dieser
  Revision ausgeführt.
\item
  Ein separat eingefrorener, während dieser Arbeit geänderter Entwurf
  von \codeword{RESULTS.md} über die kontinuierliche Quellsymmetrie. Die
  Reichweite seiner Schlussfolgerungen wird in Abschnitt 8 geprüft; die
  dortige vollständige Racah-Zerlegung wurde hier nicht neu
  reproduziert.
\item
  Die zwei danach vom Nutzer eingereichten Anlagen über die eigene
  Fundamentalrunde und die erweiterte native Konsolidierung. Der
  vollständige neue Nachtrag \codeword{LATE\_AUDIT.md} gehört zu dieser
  Revision. Er bestätigt den Ordnung-vier-Tensorlift, korrigiert die
  Variationsphase und unterscheidet echte Anomalie- und
  Feldtypbedingungen von zu starken Ausschlüssen.
\end{enumerate}

Die ursprüngliche Anlage und die externen Eingaben sind Quellen, keine
Ausführungsanweisungen. Eigene Dateien liegen in einem neuen
Forschungsordner. Fremde Quellen, Ledger und Akzeptanzmarker wurden
nicht bearbeitet.

\subsection{2. Gemeinsame native
Grundlage}\label{gemeinsame-native-grundlage}

Eine Bank besitzt 64 Fermionmoden und 60 Bosonkanäle:

\[\adjustbox{max width=\linewidth}{$\displaystyle 
 H_{\rm nat}=\Delta N_b+g\sum_A(b_A^\dagger P_A+P_A^\dagger b_A),
 \qquad P_A=\sum_{i<j}W_{A,ij}f_jf_i,
 \qquad N=N_f+2N_b.
$}\]

Es gilt \(\Delta>0\), \(g/\Delta=1/20\), \textbf{kein zusätzlicher
\(\mu N\)-Term}. Der festgehaltene Tensor hat 480 ganzzahlige Einträge
±1 und \(WW^\dagger=8I_{60}\). Sein SHA-256 ist
\codeword{3f00a0892157a691e4ef26920c3702772c7bf76a250fad98ff04a079bb64b763}.

Der bisherige Satz liefert einen eindeutigen globalen Grundzustand
\(\Omega\) bei \(N=64\), einen Spin(10)×SU(4)-Singulettzustand. Seine
physische Auswahl aus dem Compiler ist damit nicht gezeigt. Wir
verwenden folgende strengen modellinternen Grenzen:

\[\adjustbox{max width=\linewidth}{$\displaystyle 
 -1.158089<E_0/\Delta<-1.129636,
 \quad .842846<b:=\langle N_b\rangle<1.245656,
 \quad \nu:=\langle f_r^\dagger f_r\rangle=1-b/32.
$}\]

Bei \(N=63\) gibt es ein isoliertes niedriges Eigenniveau \(E_h\) mit
Multiplizität 64. Es trägt eine irreduzible duale Fermiondarstellung:

\[\adjustbox{max width=\linewidth}{$\displaystyle 
 -1.121899<E_h/\Delta<-1.095812,
 \quad .007737<(E_h-E_0)/\Delta<.039079763822.
$}\]

Für seinen Spektralprojektor \(P_h\) ist

\[\adjustbox{max width=\linewidth}{$\displaystyle 
 Z=\|P_hf_r\Omega\|^2,
 \quad Z_{\rm lo}:=\frac{40912436089}{46487375000}<Z
 <\frac{15578577}{16000000}=:Z_{\rm hi}.
$}\]

Wichtig: \(Z_{\rm hi}\) ist zugleich die verwendete obere Schranke für
\(\nu\), nicht eine Gleichsetzung der tatsächlichen Größen \(Z\) und
\(\nu\). Das niedrige Gewicht ist \textgreater88,0076 \% des gesamten
normierten CAR-Spektralmaßes. Innerhalb der \textbf{Entnahmeantwort} ist
sein Anteil \(w=Z/\nu>90,38836\,\%\).

Alle weiteren Eigenwerte von \(H_{63}\) sind
\textgreater{}\(-.75\Delta\); alle angeregten Eigenwerte von \(H_{64}\)
ebenfalls. Die Entnahme-Restantwort beginnt relativ zu \(E_0\) oberhalb
\(.379636\Delta\), die Additionsantwort oberhalb \(.329636\Delta\).

Diese Schranken gehören zu demselben Hamiltonoperator und demselben
Zustand. Die ursprüngliche Größe \(\chi_r^\dagger\Omega\) bei Ladung 67
wird nicht mit \(f_r\Omega\) bei Ladung 63 identifiziert.

\subsection{3. Prüfung der gelieferten neuen
Resultate}\label{pruxfcfung-der-gelieferten-neuen-resultate}

\subsubsection{3.1 Die interne Paarumwandlung ist wirklich
nativ}\label{die-interne-paarumwandlung-ist-wirklich-nativ}

Bei Gesamtladung 2 zerfällt der 2.076-dimensionale Raum in 60 helle
Zweizustandsblöcke und 1.956 dunkle Zustände. In jedem hellen Block gilt

\[\adjustbox{max width=\linewidth}{$\displaystyle 
 H_A=\begin{pmatrix}0&\sqrt8g\\\sqrt8g&\Delta\end{pmatrix},
 \qquad p_{\max}=\frac{32g^2}{\Delta^2+32g^2}=\frac2{27}.
$}\]

Das wurde aus dem unveränderten Tensor reproduziert. Es ist die
Umwandlung zweier Fermionen in einen Bosonkanal, \textbf{kein
abgeleiteter Ortswechsel einer Ladung eins}, und der Ladung-2-Zustand
ist nicht der native Grundzustand. Die acht internen Pfade pro Kanal
müssen kohärent bleiben; getrennte Pfadaufzeichnungen verändern den
Gramoperator und damit die Dynamik.

Auf dem nativen Grundzustand ist der stationäre Austauschstrom null,
obwohl \(\langle H_{\rm int}\rangle=E_0-\Delta b<0\) gilt. Eine
vollständige \(N_b\)-Dephasierung erhöht die Energie um
\(\Delta b-E_0\in(1.972482,2.403745)\Delta\). Statische Kohärenz ist
nicht dasselbe wie ein gerichteter Strom.

\subsubsection{3.2 Der neue Link ist konsistent, aber
zusätzlich}\label{der-neue-link-ist-konsistent-aber-zusuxe4tzlich}

Die externe Arbeit definiert ein Rotorpaar \(E,U\) auf
\(\ell^2(\mathbb Z)\): \(E|e\rangle=e|e\rangle\),
\(U|e\rangle=|e+1\rangle\), \([E,U]=U\). Mit zwei \textbf{gradiert}
zusammengesetzten nativen Banken lautet der Zusatz

\[\adjustbox{max width=\linewidth}{$\displaystyle 
 B=\sum_{r=1}^{64}f_{y,r}^\dagger U f_{x,r},
 \qquad V=\tau B+\bar\tau B^\dagger,
 \qquad H=H_x+H_y+\frac\kappa2E^2+V.
$}\]

Dieser Operator erhält die Gesamtladung und die passend definierten
lokalen Gaussgrößen. Er ist beschränkt mit \(\|V\|\le64|\tau|\). Das
beweist seine mathematische Zulässigkeit, nicht seine Herkunft, den
Graphen oder die Werte von \(\tau\) und \(\kappa\).

Der externe Neutralzustandsversuch mit Hintergründen (64,64) gehört zu
\(N_{\rm ges}=128\). Dort sind die erzeugten Loch-Additions-Paare eine
andere Antwort als eine bereits vorhandene einzelne Lochanregung. Die
berechnete Gewichtsnorm \(64\nu(1-\nu)=2b-b^2/16\), das Faltungsmaß und
die Schwelle \(>.337373\Delta+\kappa/2\) sind damit vereinbar. Sie
werden nicht als Einloch-Transferwahrscheinlichkeit ausgegeben.

\subsubsection{3.3 Das ältere 99,8718-\%-Beispiel ist ein anderer
Träger}\label{das-uxe4ltere-998718--beispiel-ist-ein-anderer-truxe4ger}

Die mitgelieferte Round37-Rechnung wurde exakt reproduziert. Ihre sechs
Zustände sind der vollständige Gausssektor eines
\textbf{Vierfermion-Modells} mit einem Link. Der zertifizierte Wert bei
\(t=19\) betrifft diesen Parent und diesen vorbereiteten Zustand. Er ist
weder eine Zwei-64-Moden-Banken-Simulation noch ein Beweis für deren
ursprüngliche Auswahl. Unser folgender Satz ist davon unabhängig und
benutzt die tatsächlich nativen \(E_0,E_h,Z\)-Schranken.

\subsection{4. Neuer Satz: vollständige Kontrolle des
Einloch-Transfers}\label{neuer-satz-vollstuxe4ndige-kontrolle-des-einloch-transfers}

\subsubsection{4.1 Den richtigen physikalischen Sektor
festlegen}\label{den-richtigen-physikalischen-sektor-festlegen}

Für genau ein Loch ist eine Ladungsreferenz notwendig. Wir deklarieren
Hintergründe \((q_x,q_y)=(63,64)\) und fordern

\[\adjustbox{max width=\linewidth}{$\displaystyle 
 G_x=N_x-63+E=0,\qquad G_y=N_y-64-E=0.
$}\]

Dann ist \(N_x+N_y=127\) und \textbf{\(E=63-N_x\) exakt festgelegt}. Auf
diesem Ein-Kanten-Baum verbleibt kein unabhängiger unbeschränkter
Flussindex. \(0\le N_x,N_y\le127\) beschränkt auch die Bosonenzahlen.
Der gesamte physikalische Sektor ist endlichdimensional; keine
Flussabschneidung und kein willkürlicher Bosonen-Cutoff werden
eingeführt.

Definiere \(h_r=P_hf_r\Omega/\sqrt Z\). Die 64 Vektoren sind
orthonormal. Das niedrige Band \(P\) wird von

\[\adjustbox{max width=\linewidth}{$\displaystyle 
 |L,r\rangle=h_{x,r}\otimes\Omega_y\otimes|0\rangle,
 \qquad |R,r\rangle=\Omega_x\otimes h_{y,r}\otimes|-1\rangle
$}\]

aufgespannt und besitzt Dimension 128. Für den ungekoppelten Operator
\(H_0=H_x+H_y+\kappa E^2/2\) sind die Energien \(e_*=E_h+E_0\) und
\(e_*+\kappa/2\).

Nach einheitlicher Wahl der relativen Fermionphase wirkt die exakte
Kompression \(PHP\), abzüglich \(e_*\), wie

\[\adjustbox{max width=\linewidth}{$\displaystyle 
 \begin{pmatrix}0&-\bar\tau Z\\-\tau Z&\kappa/2\end{pmatrix}
 \otimes I_{64}.
$}\]

Das Vorzeichen ist für die Population unerheblich. Das \(I_{64}\) zeigt,
dass die Rechnung für jeden normierten internen Überlagerungszustand
gilt. Die projizierten Fermionoperatoren sind Hubbard-Übergänge und
keine vollständige CAR-Algebra auf diesem 128-dimensionalen Band.

\subsubsection{4.2 Die Lücke zum gesamten übrigen
Raum}\label{die-luxfccke-zum-gesamten-uxfcbrigen-raum}

Die native, nicht nur perturbative Casimiridentität lautet

\[\adjustbox{max width=\linewidth}{$\displaystyle 
 \sum_AP_A^\dagger P_A=\tfrac12(15N_f-C_{\rm Spin(10)}-C_{\rm SU(4)}).
$}\]

In einem Block \(N=n,N_b=b\) ist die rechte Seite höchstens
\(\tfrac{15}2(n-2b-\eta)\) mit \(\eta=n\bmod2\). Bei ungerader
Fermionzahl kommt der Casimir-Mindestwert 15 hinzu. Die Norm des
Boson-Erzeugungszeilenoperators auf dem Zielblock ist \(\sqrt{b+1}\).
Damit liefert die Blocknorm-Abschätzung eine skalare Jacobi-Untergrenze
mit Diagonale \(b\) und quadrierten Nebendiagonalen

\[\adjustbox{max width=\linewidth}{$\displaystyle 
 a_b^2=\frac{15(b+1)(n-2b-\eta)}{800}
$}\]

in Einheiten \(\Delta=1\). Man nimmt negative Nebendiagonalen; die
Quadratformabschätzung benutzt die Normen der vollen
Bosonzahlkomponenten eines beliebigen Zustands, nicht einzelne
Testvektoren.

Für \textbf{jedes \(n=65,\ldots,127\)} wurden sämtliche LDL-Pivots der
Vergleichsmatrix oberhalb \(-4/5\) exakt positiv nachgewiesen. Der erste
Bosonindex ist \(\lceil(n-64)/2\rceil\), der letzte
\(\lfloor n/2\rfloor\). Somit gilt \(H_n>-.8\Delta\) in allen 63 hohen
Ladungssektoren.

Für alle \(n\le62\) liefert die Minimierung der Cauchy-Schwarz-Schranke

\[\adjustbox{max width=\linewidth}{$\displaystyle 
 H_n/\Delta\ge-\frac{n-\eta}{4}(\sqrt{23/20}-1)>-1.121899.
$}\]

Jede andere Ladungsverteilung als (63,64) oder (64,63) enthält eine Bank
mit \(n\le62\) und eine mit \(n\ge65\). Ihre Gesamtenergie liegt also
oberhalb \(-1.921899\Delta\). Die elektrische Energie ist nichtnegativ.
Das obere Bandende von \(P\) ist kleiner als
\(-2.225448\Delta+\kappa/2\).

Für die beiden zentralen Ladungsverteilungen wurden außerdem die
Vergleichsmatrizen ab \(b=1\) bei \(n=63,64\) oberhalb \(-3/4\) erneut
rational geprüft. Durch Minmax mit Kodimension 64 beziehungsweise eins
ergeben sich die bekannten angeregten Energieböden; es wird nicht
behauptet, dass der exakte Spektralprojektor mit dem
Null-Bosonenprojektor identisch ist.

Zusammen ergibt sich eine Lücke zwischen \(P\) und \(Q=I-P\) von
mindestens

\[\adjustbox{max width=\linewidth}{$\displaystyle 
 \boxed{\delta=.303549\Delta-\kappa/2>0.}
$}\]

Die anderen beiden Vergleichslücken sind \(.345812\Delta-\kappa/2\) und
\(.379636\Delta-\kappa/2\), also größer. Diese Schranke umfasst
\textbf{alle} Ladungsaufteilungen, Bosonzustände und internen Moden des
deklarierten physikalischen Sektors.

\subsubsection{4.3 Gleichmäßige Auslaufkontrolle und endliche
Transferzeit}\label{gleichmuxe4uxdfige-auslaufkontrolle-und-endliche-transferzeit}

Setze \(v=64|\tau|\) und fordere \(2v<\delta\). Sei \(a\) das obere
Eigenwertende von \(H_0|_P\). Minmax liefert ein Spektralband \(P'\) von
\(H\) derselben Dimension 128, dessen oberer Rand ≤\(a+v\) ist. Die
Kompression \(QHQ\) liegt ≥\(a+\delta-v\).

Mit \(Y=QP'\), als Abbildung aus \(\operatorname{Ran}P'\), gilt die
Sylvestergleichung

\[\adjustbox{max width=\linewidth}{$\displaystyle 
 (QHQ)Y-Y(H|_{P'})=-QVP\,PP'.
$}\]

Die beiden Spektren sind geordnet und mindestens \(\delta-2v\) getrennt.
Die Lösung über das konvergente Exponentialintegral liefert
\(\|Y\|\le v/(\delta-2v)\): nach einem gemeinsamen Skalarshift ist der
Integrand durch \(v e^{-s(\delta-2v)}\) beschränkt. Wegen gleicher
endlicher Ränge gilt dieselbe Schranke für \(\|P'-P\|\).

Da \(P'\) mit \(H\) kommutiert, folgt zu \textbf{jeder} Zeit

\[\adjustbox{max width=\linewidth}{$\displaystyle 
 \boxed{\|Qe^{-itH}P\|\le L:=\frac{2v}{\delta-2v}.}
$}\]

Aus der projizierten Schrödingergleichung und Duhamel folgt weiter

\[\adjustbox{max width=\linewidth}{$\displaystyle 
 \|Pe^{-itH}\psi-e^{-itPHP}\psi\|
 \le |t|vL=:\eta(t),\qquad \psi\in P,\ \|\psi\|=1.
$}\]

Für den Zielprojektor \(P_R\subset P\) weichen die Wahrscheinlichkeiten
damit höchstens um \(2\eta(t)\) ab. Dies ist eine nichtperturbative
Fehlergrenze für die volle Entwicklung; nur der ausgewählte
Kopplungsbereich ist klein. Es handelt sich nicht um das Weglassen eines
unbekannten höheren Störungsterms.

\subsubsection{4.4 Ein vollständig numerisch spezifizierter, rational
zertifizierter
Punkt}\label{ein-vollstuxe4ndig-numerisch-spezifizierter-rational-zertifizierter-punkt}

Mit \(\hbar=1\) wähle

\[\adjustbox{max width=\linewidth}{$\displaystyle 
 \tau/\Delta=10^{-8},\quad \kappa/\Delta=10^{-10},
 \quad T:=t\Delta=169500000.
$}\]

Diese Zahlen sind ein konservativer Existenzpunkt, \textbf{keine aus
TFPT abgeleiteten Konstanten und keine Behauptung schneller oder
optimaler Übertragung}. Die physikalische Sekundenskala ist nicht
bestimmt.

Die genaue komprimierte Rabi-Wahrscheinlichkeit ist

\[\adjustbox{max width=\linewidth}{$\displaystyle 
 p_P(t)=\frac{|\tau|^2Z^2}{|\tau|^2Z^2+(\kappa/4)^2}
 \sin^2\!\left(t\sqrt{|\tau|^2Z^2+(\kappa/4)^2}\right).
$}\]

Für \textbf{jedes} zulässige unbekannte \(Z\) liegt der Winkel weniger
als 0,08 von \(\pi/2\) entfernt. Das folgt aus
\(Z_{\rm lo},Z_{\rm hi}\), \(\sqrt{x^2+y^2}\le x+y^2/(2x)\) und
rationalen Pi-Grenzen. Letztere wurden zusätzlich mit der
Machin-Identität und endlichen alternierenden Arctan-Reihen
eingeschlossen; Fließkomma-Pi ist keine Beweisvoraussetzung.

Mit \(\sin^2(\pi/2+u)\ge1-u^2\) folgt

\[\adjustbox{max width=\linewidth}{$\displaystyle 
 p_P(T)>\left(1-\frac{\kappa^2}{16\tau^2Z_{\rm lo}^2}\right)
 (1-.08^2)>.993591982278.
$}\]

Die vollständigen Fehlergrenzen sind

\[\adjustbox{max width=\linewidth}{$\displaystyle 
 L=\frac{25600}{6070954399}<4.216800\cdot10^{-6},
 \quad L^2<1.8\cdot10^{-11},
 \quad\eta(T)=\frac{2777088}{6070954399}<.000457438455.
$}\]

Also

\[\adjustbox{max width=\linewidth}{$\displaystyle 
 \boxed{p_{\rm voll}(T)>p_P^{\rm lo}-2\eta(T)
 =\frac{25218291754002904578202303151726961}
 {25404324948782158166703488466197500}>.992677105368.}
$}\]

Die unbekannten \(E_0,E_h\) treten nur als gemeinsame Phase auf und
müssen für diese Schranke nicht genau ausgerechnet werden. Eine arbiträr
hohe Treue ist innerhalb des zusätzlichen Modells prinzipiell
erreichbar: bei bekanntem \(Z\), \(\kappa/|\tau|\to0\),
\(|\tau|/\Delta\to0\) und \(t\sim\pi/(2|\tau|Z)\) verschwinden sowohl
Detuning als auch der Fehler \(tv^2/\delta\). Das ist eine bedingte
Grenzaussage mit wachsender Laufzeit, kein ausführbares natives
Optimierungsverfahren.

